\scanpage{188}

\noindent\begin{tabular}{@{}|p{.92\linewidth}@{}}les classes $E_n$ des modules $k[\unsure{x|X}]/x^n$, $n\unsure{:|=}1,2,\ldots$. L'application $\unsure{v|\vee}$ dans $A$ est réduite à l'identité. On a $d(E_n)=n$, $E_1$ est l'unité.\end{tabular}

Structure multiplicative. Posons $\xi=E_2$, classe de $k[x]/(X^{\unsure{2}})$. Alors $\xi^n$ est la classe de $F_n=k[X_1,\ldots,X_n]/(X_1^2,\ldots X_n^2)$, muni de l'opérateur nilpotent $T$ de multiplication par la classe $\dot S$ de $S=X_1+\ldots+X_n$. \deleted{Les} Le \unsure{$n$-ième} itéré \inserted{$T^h$} de $T$ est la multiplication par

\[ \deleted{\dot S^h=h!\sum_{\substack{\unsure{H\subset[1,n]}\\\unsure{\lvert H\rvert=h|\#H}}}}\qquad \dot S^h=k!\sum_{1\leq i_1<\cdots<i_{\unsure{k|h}}\leq n}\dot X_{i_1}\ldots\dot X_{i_{\unsure{h|k}}}. \]

En particulier $T^n$ est la multiplication par $\dot S^n=n!\dot X_1\ldots\dot X_n$, et $T^{n+1}$ est nul.

On en conclut aussitôt que $T^n F_n$ est de dim 1 (base formée de $\dot X_1\dot X_2\ldots\dot X_n$) et $T^{n+1}(F_n)=0$ ; d'où résulte facilement que \deleted{\unsure{la multiplicité}} $\xi^n$ est combinaison linéaire des $E_1,\ldots E_n$ à coefficients $\in Z^+$, celui de $E_n$ étant 1 :

\margin{\unsure{3}}

% archive sheet number: 187
